\documentclass[10pt,a4paper]{amsart}
\usepackage[margin=19mm]{geometry}
\usepackage{amsmath,amssymb,mathtools}
\usepackage[scr=rsfs,cal=euler]{mathalfa}
\usepackage{xfrac}
\usepackage[unicode,hidelinks,bookmarks=false]{hyperref}
\newcommand{\ie}{i.e.\ }
\newcommand{\cf}{cf.\ }
\newcommand{\resp}{respectively\ }
\newcommand{\Subsection}{\subsection}
\providecommand{\nobreakdash}{\nobreak\hbox{-}}
\setlength{\emergencystretch}{2em}
\multlinegap0pt
\usepackage{amssymb}
\newtheorem{corollary}{Corollary}
\newcommand{\dd}{\,d}
\newcommand{\dive}{\nabla\!\cdot}
\title{REFLECTED DIFFUSIONS,\\ NO-FLUX CONTINUITY EQUATIONS AND\\ CONFINED LAGRANGIAN FLOWS IN BOUNDED DOMAINS}
\author{Rama CONT}
\date{Source: arXiv:2607.28344v1 (CC BY 4.0)}
\begin{document}
\maketitle

\noindent\emph{Source and licence.} REFLECTED DIFFUSIONS,\\ NO-FLUX CONTINUITY EQUATIONS AND\\ CONFINED LAGRANGIAN FLOWS IN BOUNDED DOMAINS. Version arXiv:2607.28344v1 (CC BY 4.0), distributed by arXiv under the Creative Commons Attribution 4.0 International licence, http://creativecommons.org/licenses/by/4.0/, as recorded in the arXiv licence metadata for this exact version and shown on the abstract page https://arxiv.org/abs/2607.28344v1. Full licence notice: \texttt{licenses/arxiv-2607.28344-CC-BY-4.0.txt}.\par\smallskip

\noindent\textit{Editorial scope.} Counted unit: the article's corollary cor:entangled (source0.tex lines 2101--2117). The article's complete original proof of it is delivered with it (source0.tex lines 2119--2217, one \texttt{proof} environment of 99 lines). No mathematics is written, repaired or re-derived: the statement and the proof are the article's own lines, in the article's own notation, and no retained line is substituted or re-flowed.  No further source-local context is needed: every cross-reference of every retained line resolves inside the two delivered spans.  Nothing was omitted from any retained span, so the package carries no omission record.  No retained line contains a citation, so the wrapper carries no bibliography and no bibliography entry.  No retained line contains a figure environment, an image-inclusion command or drawing code, so no image is delivered and the package declares no asset dependency.  Editorial formatting only, in addition: an AMS \texttt{amsart} preamble that restates the article's own declarations and macros used by the retained lines, namely the environments \texttt{corollary} on separate counters, together with the standard packages those lines need.  The retained environments print in this document as Corollary 1 in order of appearance, whereas the article prints the same items with the numbering of its own full document.  The complete native source of the article as distributed in the arXiv e-print for this exact version is bundled unchanged as \texttt{sources/206409/source0.tex}.
\begin{corollary}[Failure of tangency forces non-integrability]
\label{cor:entangled}
Let $\Gamma\subset\partial\Omega$ be a relatively open $C^{1,1}$ patch,
let $p$ and $J$ be of class $C^1$ in space and time up to $\Gamma$
on a time interval, with $J\cdot\nu=0$ on $\Gamma$, and suppose $p$
vanishes on $\Gamma$ nondegenerately:
$p(x)\ge c\operatorname{dist}(x,\Gamma)$ with $c>0$. Let
$x_0\in\Gamma$ and $\nu_0:=\nu(x_0)$.
Fix $t_0$ 
 in the stated time interval. If
\[
v(x,t_0)\cdot\nu_0\ \not\longrightarrow\ 0
\qquad\text{as }\Omega\ni x\to x_0
\]
\emph{(}the approach being unrestricted\emph{)}, then $v(.,t_0)\notin L^1$
on any spatial collar of $x_0$.
\end{corollary}

\begin{proof}
We argue invariantly. Fix $x_0\in\Gamma$, write $\nu_0:=\nu(x_0)$ and
let $\{\tau_1,\dots,\tau_{d-1}\}$ be an orthonormal basis of the
tangent space $\nu_0^\perp$; only the \emph{single} vector $\nu_0$
is ever differentiated against.
No derivative of the normal field is required; the $C^{1,1}$
 assumption is used only to obtain the tubular nearest-point projection.
The curved boundary enters only through tangential differentiation
of functions defined on $\bar\Omega$, which for a $C^1$ patch is
well defined without any chart: for $\tau\in\nu_0^\perp$ we set
$\partial_\tau w(x_0):=(w\circ\gamma)'(0)$ for any $C^1$ curve
$\gamma$ in $\Gamma$ with $\gamma(0)=x_0$, $\gamma'(0)=\tau$, which
agrees with the ambient derivative $\tau\cdot\nabla w(x_0)$ whenever
$w$ is $C^1$ up to $\Gamma$ and is therefore independent of the
curve. This is the only property of $\Gamma$ used below, and it is
why no flattening --- which would not preserve the divergence, and
would require second-order regularity of $\partial\Omega$ to
control the error --- is needed.

We also record, for later use, that the $C^1$ regularity of $p$
together with $p|_\Gamma\equiv0$ gives the \emph{upper} bound
$p(x)\le\|\nabla p\|_\infty\operatorname{dist}(x,\partial\Omega)$
near $x_0$, so that the nondegeneracy hypothesis makes
$p\asymp\operatorname{dist}$ there.

\emph{Step 1: $\dive J=0$ on $\Gamma$.} Since $p|_\Gamma\equiv0$
throughout a time interval, $\partial_tp=0$ on $\Gamma$; the
continuity equation holds in $\Omega$ and both of its terms extend
continuously to $\Gamma$, so $\dive J=-\partial_tp=0$ there.

\emph{Step 2:} One of the following holds:
\begin{itemize}
\item[(A)] $J$ vanishes identically on some relatively open
$\Gamma'\subset\Gamma$ containing $x_0$; or
\item[(B)] every relative neighbourhood of $x_0$ in $\Gamma$
contains a point at which $J\neq0$.
\end{itemize}
We show that \textup{(A)} implies $v(x)\cdot\nu_0\to0$ as
$\Omega\ni x\to x_0$, and that \textup{(B)} implies $v\notin L^1$ on
every collar of $x_0$. The corollary follows: if the limit fails,
\textup{(A)} is excluded by the first implication, so \textup{(B)}
holds and the second applies. Arguing through the dichotomy, rather
than through the value of a directional limit, is what allows the
hypothesis to be the unrestricted limit; a computation along the
inward normal ray alone would prove only the weaker statement in
which that particular limit is assumed nonzero.

\emph{Step 3: \textup{(A)} implies tangency at $x_0$.} Let $J\equiv0$
on $\Gamma'$. For $x'\in\Gamma'$ and $\tau$ tangent at $x'$,
differentiating along a $C^1$ curve in $\Gamma'$ gives
$\partial_\tau J(x')=0$; hence, with an orthonormal tangent frame
$\{\tau_i'\}$ at $x'$ and $\nu':=\nu(x')$, Step 1 yields
\begin{equation}\label{eq:entangledlimit}
\nu'^{\top}(\nabla J)(x')\,\nu'
=\operatorname{tr}(\nabla J)(x')-\sum_{i=1}^{d-1}
\tau_i'\cdot\big(\partial_{\tau_i'}J\big)(x')=0 ,
\end{equation}
using $\tau^{\top}(\nabla J)\tau=\tau\cdot\partial_\tau J$. No
curvature term appears here: the second fundamental form would enter
multiplied by $J\cdot\nu$, which vanishes on $\Gamma$. Now let $\pi$
be a nearest-point projection onto $\partial\Omega$, defined for $x$
near $x_0$, so that $\pi(x)\in\Gamma'$ and
$x-\pi(x)=-d(x)\,\nu(\pi(x))$ with
$d(x):=\operatorname{dist}(x,\partial\Omega)$. Since $\nabla J$ is
uniformly continuous near $x_0$ and $J(\pi(x))=0$,
\[
J(x)=-d(x)\,(\nabla J)(\pi(x))\,\nu(\pi(x))+o\big(d(x)\big),
\]
hence
\[
J(x)\cdot\nu_0=-d(x)\,\nu_0^{\top}(\nabla J)(\pi(x))\,\nu(\pi(x))
+o\big(d(x)\big).
\]
As $x\to x_0$ we have $\pi(x)\to x_0$ and, $\Gamma$ being $C^1$,
$\nu(\pi(x))\to\nu_0$; by continuity of $\nabla J$ and
\eqref{eq:entangledlimit} at $x'=x_0$, the bracketed factor tends to
$\nu_0^{\top}(\nabla J)(x_0)\nu_0=0$. Therefore $J(x)\cdot\nu_0
=o(d(x))$, and the nondegeneracy $p\ge c\,d$ gives
\[
\big|v(x)\cdot\nu_0\big|=\frac{|J(x)\cdot\nu_0|}{p(x)}
\le\frac{o(d(x))}{c\,d(x)}\longrightarrow0 ,
\]
the approach being unrestricted.

\emph{Step 4: \textup{(B)} implies non-integrability.} Fix any collar
of $x_0$ and choose, inside it, a point $x_1\in\Gamma$ with
$J(x_1)\neq0$. By continuity there are $\eta>0$ and $r>0$ with
$|J|\ge\eta$ on $B_r(x_1)\cap\Omega$, and by the upper bound
recorded above $p\le\|\nabla p\|_\infty d$ there, whence
\[
|v|=\frac{|J|}{p}\ \ge\ \frac{\eta}{\|\nabla p\|_\infty\,d(\cdot)}
\qquad\text{on }B_r(x_1)\cap\Omega .
\]
Integrating in the normal direction, $\int_0^r s^{-1}\dd s=\infty$,
so $v\notin L^1(B_r(x_1)\cap\Omega)$ and a fortiori $v\notin L^1$ on
the collar. As the collar was arbitrary, this holds for every collar
of $x_0$.

\end{proof}

\end{document}
